\documentclass[../midgard.tex]{subfiles}
\graphicspath{{\subfix{../images/}}}
\begin{document}

\section*{Introduction}
\label{h:introduction}
\addcontentsline{toc}{chapter}{Introduction}

Midgard is a Layer 2 (L2) scaling solution for the Cardano blockchain.
It employs optimistic rollup technology to enhance Cardano's capacity to process transactions and host more complex applications, delivering a richer user experience at a more competitive cost.
As Cardano continues to grow in usage and demand, scaling solutions like Midgard are critical for maintaining high performance and low transaction costs.
This whitepaper describes the architecture and technical design of the Midgard protocol, detailing how it integrates with Cardano's Layer 1 (L1) to deliver secure and efficient transaction processing.

Optimistic rollups process blocks of transactions off-chain and commit those blocks' headers onchain to the L1 ledger.
Each block is committed by a Midgard operator who guarantees the block's validity.
The block then waits in a queue for at least a fixed duration to be merged into the confirmed state of the optimistic rollup on L1.
The design requires the operator to collateralize its claim with a bond deposit and make the full block contents permissionlessly retrievable through the data availability (DA) layer for at least the challenge and recovery horizon.

While a committed block is queued, anyone can inspect its contents on the data availability layer and ascertain whether it is valid.
If someone detects that the block is invalid, they can submit a fraud proof to prevent it from being merged into the confirmed state, slash the operator's bond, and receive a portion of the forfeited bond as a reward.
Thus, an optimistic rollup can process a large number of transactions offchain while maintaining security and finality properties that are similar to transactions processed directly onchain, as long as:
\begin{itemize}
    \item The bond requirement for the rollup's operators is large enough to discourage fraud.
    \item The reward for preventing an invalid block from merging is large enough to encourage public vigilance in watching the operators.
    \item The waiting period for committed blocks is long enough to allow the watchers to detect and prove fraud before those blocks are merged.
    \item The data availability layer is accessible by anyone who wishes to inspect the rollup blocks at any time that they wish to do so.
\end{itemize}

The economic deterrent holds only when invalid data remains obtainable, at least one independent watcher detects the violation, an onchain proof path covers it and completes before maturity, and the expected loss exceeds the maximum extractable benefit.
A disqualified block does not affect the confirmed state, but that fact alone does not establish sufficient deterrence; bond and reward values must account for detection probability, proof and capital costs, correlated failures, and the value at risk.

The main design goal of Midgard is to streamline the processes by which blocks are committed/merged, fraud is detected, and fraud proofs are verified onchain.
Advancing this goal allows the security parameters to be calibrated to achieve a better balance between security, transaction throughput, confirmation time, and community participation in committing blocks and detecting fraud.

\subsection*{Scalability and efficiency}
\label{h:scalability-and-efficiency}

By processing transactions off-chain and only validating them on-chain when fraud proofs challenge them, Midgard significantly increases throughput and reduces costs for Cardano transactions.
Its rollup blocks use sparse Merkle trees and compact state representations to enhance the protocol's efficiency further, enabling it to handle a large volume of transactions without overburdening the L1.

Deterministic transaction encoding, typed roots, and authenticated state witnesses allow many violations to be localized to a bounded claim and Merkle path.
Some violation families still depend on transaction-wide material, neighbouring transition steps, prior committed state, or authenticated L1 events; each proof procedure must state its complete witness and binding requirements.
The intended result is bounded onchain verification rather than re-execution of a complete block, but proof size and cost are deployment claims that require budget and end-to-end evidence for every enabled family.

\subsection*{Censorship resistance and fallback mechanisms}

On its own, the optimistic rollup mechanism described above ensures a high-level of assurance for the validity of block headers committed to the state queue and merged to Midgard's confirmed state.
However, it does not prevent operators from censoring users' deposits, withdrawals, and L2 transactions.
Consequently, Midgard's consensus protocol includes additional smart contract mechanisms to provide censorship resistance for these events.

Midgard deposits and withdrawals are initiated via L1 smart contracts that assign definite inclusion times to them.
An operator block is invalid if it contains these inclusion times in its event interval but fails to include the associated deposit or withdrawal events.
This ensures that if operators continue committing blocks to Midgard's state queue, then they cannot ignore deposit and withdrawal events.

Midgard L2 transaction requests are typically submitted to operators via a publicly accessible API, and they can be ignored by operators.
However, any user can escalate his L2 transaction request by posting a transaction order on L1.
Similar to Midgard deposits and withdrawals, an L1 transaction order is assigned an inclusion time that makes omission from a covering block provable while valid block production continues.

If Midgard operators stop committing blocks at all to the state queue, then the inclusion times on their own cannot guarantee that deposits, withdrawals, and L2 transactions will be processed in a timely manner.
However, for this extreme case, Midgard's consensus protocol includes the escape hatch mechanism, which temporarily lowers the bond requirement for new operators.
The target escape-hatch design makes it easier for honest actors to become Midgard operators and restart block production.
It is not a direct user force-exit and does not, on its own, guarantee timely recovery: production safety requires deployed validators and transactions, an economically viable replacement operator, available data and keys, and a tested recovery runbook.

\subsection*{Feedback and contributions welcome!}
\label{h:feedback-and-contributions-welcome}

This document specifies the ledger, onchain contracts, and offchain components of the Midgard architecture.
We detail the roles of users, operators, and watchers, how blocks are committed and merged on L1, and how fraud proofs are handled.

We invite developers, validators, and the broader community to engage with this document and offer feedback and contributions to help refine and enhance the protocol.

\end{document}
